\clearpage
\refstepcounter{chapter}
\addcontentsline{toc}{chapter}{\texorpdfstring{\protect\tocmark{cyanaccent}}{}Bibliography}
\renewcommand{\chaptertitle}{Bibliography}
\renewcommand{\sectiontitle}{}
\pagestyle{plain}
\noindent{\Huge\bfseries Bibliography\par}
\vspace{0.5em}

\begingroup
\small
\raggedright
\interlinepenalty=10000 % Keep each reference together across page breaks.
\setlength{\parskip}{5pt}
\newcommand{\srchead}[1]{\par\addvspace{10pt}\noindent{\normalsize\bfseries #1}\par\nobreak\vspace{2pt}\nobreak}

\srchead{Regression}
Hyndman, R.~J. and Koehler, A.~B. (2006). Another look at measures of forecast accuracy. \emph{International
Journal of Forecasting} 22(4).\par
Willmott, C.~J. and Matsuura, K. (2005). Advantages of the mean absolute error (MAE) over the root mean
square error (RMSE) in assessing average model performance. \emph{Climate Research} 30.\par
Makridakis, S. (1993). Accuracy measures: theoretical and practical concerns. \emph{International Journal of
Forecasting} 9(4).\par
Hyndman, R.~J. and Athanasopoulos, G. \emph{Forecasting: Principles and Practice}, chapter on evaluating forecast
accuracy, \url{https://otexts.com/fpp3/accuracy.html}.\par
On sMAPE's asymmetry: \url{https://scholarworks.utep.edu/cgi/viewcontent.cgi?article=1865&context=cs_techrep}.
On wMAPE against MAPE: \url{https://www.baeldung.com/cs/mape-vs-wape-vs-wmape}.\par
scikit-learn, \emph{Metrics and scoring: regression metrics},
\url{https://scikit-learn.org/stable/modules/model_evaluation.html}.\par
McCullagh, P. and Nelder, J.~A. (1989). \emph{Generalized Linear Models}, 2nd ed. Chapman \& Hall; the unit deviances
of the Poisson and Gamma families.\par
Koenker, R. and Bassett, G. (1978). Regression quantiles. \emph{Econometrica} 46(1).\par
Gneiting, T. (2011). Quantiles as
optimal point forecasts. \emph{International Journal of Forecasting} 27(2).\par
Blaskowitz, O. and Herwartz, H. (2011). On economic evaluation of directional forecasts. \emph{International Journal of
Forecasting} 27(4).\par

\srchead{Classification}
Fawcett, T. (2006). An introduction to ROC analysis. \emph{Pattern
Recognition Letters} 27(8).\par
Hanley, J.~A. and McNeil, B.~J. (1982). The meaning and use of the area under a receiver operating
characteristic (ROC) curve. \emph{Radiology} 143(1).\par
van Rijsbergen, C.~J. (1979). \emph{Information Retrieval}, 2nd ed. Butterworths.\par
Powers, D.~M.~W.
(2011). Evaluation: from precision, recall and F-measure to ROC, informedness, markedness and correlation. \emph{Journal of Machine
Learning Technologies} 2(1).\par
Davis, J. and Goadrich, M. (2006). The relationship between precision-recall and ROC curves. \emph{ICML}.\par
Saito, T. and
Rehmsmeier, M. (2015). The precision-recall plot is more informative than the ROC plot when evaluating binary classifiers on imbalanced
datasets. \emph{PLOS ONE} 10(3).\par
Brier, G.~W. (1950). Verification of forecasts expressed in terms of probability. \emph{Monthly Weather
Review} 78(1).\par
Good, I.~J. (1952). Rational decisions. \emph{Journal of the Royal Statistical Society B} 14(1).\par
Gneiting, T. and
Raftery, A.~E. (2007). Strictly proper scoring rules, prediction, and estimation. \emph{JASA} 102(477).\par
Brodersen, K.~H. et al. (2010). The balanced accuracy and its posterior distribution. \emph{ICPR}.\par
Jaccard, P. (1901). Étude comparative de la distribution florale dans une portion des Alpes et des Jura.
\emph{Bulletin de la Société Vaudoise des Sciences Naturelles} 37.\par
scikit-learn, \texttt{d2\_log\_loss\_score}, \url{https://scikit-learn.org/stable/modules/model_evaluation.html}.\par
Sitarz, M. (2023). Extending F1 metric, probabilistic approach. arXiv:2210.11997.\par
Cohen, J. (1960). A coefficient of agreement for nominal scales. \emph{Educational and Psychological Measurement}
20(1).\par
Matthews, B.~W. (1975). Comparison of the predicted and observed secondary structure of T4 phage lysozyme. \emph{Biochimica et
Biophysica Acta} 405(2).\par
Chicco, D. and Jurman, G. (2020). The advantages of the Matthews correlation coefficient (MCC) over F1 score
and accuracy in binary classification evaluation. \emph{BMC Genomics} 21(6).\par
Ferrer, L. (2022). Analysis and comparison of classification metrics. arXiv:2209.05355.\par

\srchead{Clustering}
Vinh, N.~X., Epps, J. and Bailey, J. (2010). Information theoretic measures for clusterings comparison: variants,
properties, normalization and correction for chance. \emph{JMLR} 11.\par
Rand, W.~M. (1971). Objective criteria for the evaluation of clustering methods. \emph{JASA}
66(336).\par
Hubert, L. and Arabie, P. (1985). Comparing partitions. \emph{Journal of Classification} 2(1).\par
Caliński, T. and Harabasz, J. (1974). A dendrite method for cluster analysis. \emph{Communications in Statistics} 3(1).\par
Rosenberg, A. and Hirschberg, J. (2007). V-measure: a conditional entropy-based external
cluster evaluation measure. \emph{EMNLP-CoNLL}.\par
Davies, D.~L. and Bouldin, D.~W. (1979). A cluster separation measure. \emph{IEEE TPAMI} 1(2).\par
Fowlkes, E.~B. and Mallows, C.~L. (1983). A method for comparing two hierarchical clusterings. \emph{JASA} 78(383).\par
Rousseeuw, P.~J. (1987). Silhouettes: a graphical aid to the interpretation and validation of cluster analysis.
\emph{Journal of Computational and Applied Mathematics} 20.\par
scikit-learn, \emph{Clustering performance evaluation},
\url{https://scikit-learn.org/stable/modules/clustering.html#clustering-performance-evaluation}.\par

\srchead{Ranking}
Manning, C.~D., Raghavan, P. and Schütze, H. (2008). \emph{Introduction to Information Retrieval}, chapter 8.
Cambridge University Press, \url{https://nlp.stanford.edu/IR-book/}.\par
Voorhees, E.~M. (1999). The TREC-8 question answering track
report.\par
On MAP: \url{https://juneandrews.com/2014/12/15/mean-average-precision-isnt-so-nice/} and
\url{https://sdsawtelle.github.io/blog/output/mean-average-precision-MAP-for-recommender-systems.html}.\par
Herlocker, J.~L. et al. (2004). Evaluating collaborative filtering recommender systems. \emph{ACM
TOIS} 22(1).\par
Ge, M., Delgado-Battenfeld, C. and Jannach, D. (2010). Beyond accuracy: evaluating recommender systems by coverage and
serendipity. \emph{RecSys}.\par
Järvelin, K. and Kekäläinen, J. (2002). Cumulated gain-based evaluation of IR techniques. \emph{ACM TOIS} 20(4).
\url{https://aman.ai/recsys/metrics/#normalized-discounted-cumulative-gain-ndcg-1}.\par
Koren, Y. and Sill, J. (2011). OrdRec: an ordinal model for predicting personalized item rating distributions. \emph{RecSys}.
\url{https://www.ijcai.org/Proceedings/13/Papers/449.pdf}; \url{https://aman.ai/recsys/metrics/#fraction-of-concordant-pairs-fcp}.\par
Castells, P., Hurley, N. and Vargas, S. (2015). Novelty and diversity in recommender systems. In
\emph{Recommender Systems Handbook}, 2nd ed. Springer.\par

\srchead{Computer Vision}
Long, J., Shelhamer, E. and Darrell, T. (2015). Fully convolutional networks for semantic segmentation.
\emph{CVPR}.\par
Everingham, M. et al. (2010). The PASCAL Visual Object Classes (VOC) challenge. \emph{IJCV} 88(2).\par
Yang, Y. and Ramanan, D. (2013). Articulated human detection with flexible mixtures of parts. \emph{IEEE TPAMI} 35(12).\par
Andriluka, M. et al. (2014). 2D human pose estimation: new benchmark and state of the art analysis. \emph{CVPR} (PCKh).\par
Lin, T.-Y. et al. (2014). Microsoft COCO: common objects in context. \emph{ECCV}. COCO keypoint evaluation,
\url{https://cocodataset.org/#keypoints-eval}.\par
Wang, Z., Bovik, A.~C., Sheikh, H.~R. and Simoncelli, E.~P. (2004). Image quality assessment: from error visibility to
structural similarity. \emph{IEEE Transactions on Image Processing} 13(4), \url{https://www.cns.nyu.edu/pub/lcv/wang03-preprint.pdf}.\par
Huynh-Thu, Q. and Ghanbari, M. (2008). Scope of validity of PSNR in image/video quality assessment. \emph{Electronics Letters} 44(13).
Limits of SSIM: \url{https://espace2.etsmtl.ca/id/eprint/12070/1/Noumeir%20R.%202015%2012070%20Limitations%20of%20the%20SSIM%20quality%20metric.pdf};
worked example: \url{https://scikit-image.org/docs/0.25.x/auto_examples/transform/plot_ssim.html}.\par
Dice, L.~R. (1945). Measures of the amount of ecologic association between species. \emph{Ecology}
26(3).\par
Sørensen, T. (1948). A method of establishing groups of equal amplitude in plant sociology. \emph{Biologiske Skrifter} 5(4).
\url{https://www.nature.com/articles/s41598-023-35605-7}.\par
Kirillov, A., He, K., Girshick, R., Rother, C. and Dollár, P. (2019). Panoptic segmentation. \emph{CVPR},
\url{https://openaccess.thecvf.com/content_CVPR_2019/papers/Kirillov_Panoptic_Segmentation_CVPR_2019_paper.pdf}.\par

\srchead{NLP}
Papineni, K., Roukos, S., Ward, T. and Zhu, W.-J. (2002). BLEU: a method for automatic evaluation of machine translation.
\emph{ACL}, \url{https://aclanthology.org/P02-1040.pdf}.\par
Banerjee, S. and Lavie, A. (2005). METEOR: an automatic metric for MT evaluation with improved correlation with human
judgments. \emph{ACL Workshop}, \url{https://aclanthology.org/W05-0909.pdf}.\par
Lin, C.-Y. (2004). ROUGE: a package for automatic evaluation of summaries. \emph{ACL Workshop},
\url{https://aclanthology.org/W04-1013.pdf}.\par
Snover, M., Dorr, B., Schwartz, R., Micciulla, L. and Makhoul, J. (2006). A study of translation edit rate with targeted
human annotation. \emph{AMTA}, \url{https://aclanthology.org/2006.amta-papers.25.pdf}.\par
Rajpurkar, P., Zhang, J., Lopyrev, K. and Liang, P. (2016). SQuAD: 100,000+ questions for machine comprehension of text.
\emph{EMNLP}, \url{https://nlp.stanford.edu/pubs/rajpurkar2016squad.pdf}.\par
Levenshtein, V.~I. (1966). Binary codes capable of correcting deletions, insertions, and reversals. \emph{Soviet Physics
Doklady} 10(8).\par
Morris, A.~C., Maier, V. and Green, P. (2004). From WER and RIL to MER and WIL. \emph{Interspeech}.\par
\url{https://aclanthology.org/E17-2007.pdf}.\par

\srchead{GenAI}
Jelinek, F., Mercer, R.~L., Bahl, L.~R. and Baker, J.~K. (1977). Perplexity, a measure of the difficulty of speech
recognition tasks. \emph{JASA} 62(S1).\par
Holtzman, A. et al. (2020). The curious case of neural text degeneration. \emph{ICLR},
arXiv:1904.09751. \url{https://aclanthology.org/2021.deelio-1.5.pdf}; \url{https://huggingface.co/docs/transformers/perplexity};
\url{https://kilthub.cmu.edu/articles/journal_contribution/Evaluation_Metrics_For_Language_Models/6605324};
\url{https://brenocon.com/blog/2013/01/perplexity-as-branching-factor-as-shannon-diversity-index/}.\par
Zhang, T., Kishore, V., Wu, F., Weinberger, K.~Q. and Artzi, Y. (2020). BERTScore: evaluating text generation with
BERT. \emph{ICLR}, arXiv:1904.09675. Reference implementation: \url{https://github.com/Tiiiger/bert_score}.\par
Pillutla, K. et al. (2021). MAUVE: measuring the gap between neural text and human text using divergence frontiers.
\emph{NeurIPS}, arXiv:2102.01454.\par
Pillutla, K. et al. (2023). MAUVE scores for generative models: theory and practice.
arXiv:2212.14578. \url{https://krishnap25.github.io/mauve/}.\par
Salimans, T. et al. (2016). Improved techniques for training GANs. \emph{NeurIPS}, arXiv:1606.03498.\par
Barratt, S. and Sharma,
R. (2018). A note on the Inception Score. arXiv:1801.01973.\par
Heusel, M. et al. (2017). GANs trained by a two time-scale update rule converge to a local Nash equilibrium. \emph{NeurIPS}.\par
Parmar, G., Zhang, R. and Zhu, J.-Y. (2022). On aliased resizing and surprising subtleties in GAN evaluation. \emph{CVPR},
arXiv:2104.11222.\par
Chong, M.~J. and Forsyth, D. (2020). Effectively unbiased FID and Inception Score. \emph{CVPR}, arXiv:1911.07023.\par
Jayasumana, S. et al. (2024). Rethinking FID. \emph{CVPR}, arXiv:2401.09603.\par
Zhang, R., Isola, P., Efros, A.~A., Shechtman, E. and Wang, O. (2018). The unreasonable effectiveness of deep features as a
perceptual metric. \emph{CVPR}, arXiv:1801.03924.\par
Hessel, J. et al. (2021). CLIPScore: a reference-free evaluation metric for image captioning. \emph{EMNLP},
arXiv:2104.08718. \url{https://lightning.ai/docs/torchmetrics/stable/multimodal/clip_score.html}.\par
Cho, J. et al. (2024). Davidsonian scene graph: improving reliability in fine-grained evaluation for text-to-image
generation. \emph{ICLR}, arXiv:2310.18235.\par
Lin, Z. et al. (2024). Evaluating text-to-visual generation with image-to-text generation. \emph{ECCV},
arXiv:2404.01291. \url{https://blog.ml.cmu.edu/2024/10/07/vqascore-evaluating-and-improving-vision-language-generative-models/}.\par
ITU-T Recommendation P.800 (1996), \emph{Methods for subjective determination of transmission quality}; P.808 (2018),
\emph{Subjective evaluation of speech quality with a crowdsourcing approach}. arXiv:2204.02249.\par
ITU-T Recommendation P.862 (2001), \emph{Perceptual evaluation of speech quality}; P.862.1 (mapping to MOS-LQO);
P.862.2 (wideband). \url{https://lightning.ai/docs/torchmetrics/stable/audio/perceptual_evaluation_speech_quality.html}.\par
Taal, C.~H., Hendriks, R.~C., Heusdens, R. and Jensen, J. (2011). An algorithm for intelligibility prediction of
time-frequency weighted noisy speech. \emph{IEEE TASLP} 19(7).\par
Shannon, R.~V. et al. (1995). Speech recognition with primarily temporal
cues. \emph{Science} 270(5234).
\url{https://lightning.ai/docs/torchmetrics/stable/audio/short_time_objective_intelligibility.html};
\url{https://www.mathworks.com/help/audio/ug/measure-speech-intelligibility-and-perceived-audio-quality-with-stoi-and-visqol.html}.\par
GLUE arXiv:1804.07461 and \url{https://gluebenchmark.com/}; SuperGLUE arXiv:1905.00537 and
\url{https://super.gluebenchmark.com/}; IFEval arXiv:2311.07911; BIG-Bench arXiv:2206.04615; BBH arXiv:2210.09261; BBEH
arXiv:2502.19187; MATH arXiv:2103.03874; GPQA arXiv:2311.12022; MuSR arXiv:2310.16049; MMLU-Pro arXiv:2406.01574 and
\url{https://huggingface.co/spaces/TIGER-Lab/MMLU-Pro}; SQuAD 2.0 arXiv:1806.03822 and \url{https://rajpurkar.github.io/SQuAD-explorer/};
ARC arXiv:1911.01547.\par
DeepSeek-AI (2025). DeepSeek-R1: incentivizing reasoning capability in LLMs via reinforcement learning. arXiv:2501.12948.
Official results: \url{https://github.com/deepseek-ai/DeepSeek-R1}.\par
ARC Prize Foundation (2024). ARC Prize 2024: technical report. arXiv:2412.04604.\par
ARC Prize Foundation (2025). ARC-AGI-2, \url{https://arcprize.org/arc-agi/2}.\par
ARC Prize Foundation (2026). ARC-AGI-3, \url{https://arcprize.org/arc-agi/3}.\par

\srchead{Probabilistic}
Naeini, M.~P., Cooper, G.~F. and Hauskrecht, M. (2015). Obtaining well calibrated probabilities using Bayesian binning.
\emph{AAAI}.\par
Guo, C., Pleiss, G., Sun, Y. and Weinberger, K.~Q. (2017). On calibration of modern neural networks. \emph{ICML}.
\url{https://proceedings.mlr.press/v70/guo17a.html}.\par
Nixon, J. et al. (2019; revised 2020). Measuring calibration in deep learning. arXiv:1904.01685v2.
\url{https://arxiv.org/abs/1904.01685v2}.\par
DeGroot, M.~H. and Fienberg, S.~E. (1983). The comparison and evaluation of forecasters. \emph{The Statistician} 32(1/2), 12--22.
\url{https://doi.org/10.2307/2987588}.\par
Scikit-learn documentation, \emph{Probability calibration}.
\url{https://scikit-learn.org/stable/modules/calibration.html}.\par
NannyML documentation, \emph{How it works: performance estimation},
\url{https://nannyml.readthedocs.io/en/stable/how_it_works/performance_estimation.html}.\par
Kivimäki, J., Białek, J., Kuberski, W. and Nurminen, J.~K. (2025). Performance estimation in binary classification using
calibrated confidence. arXiv:2505.05295.
\url{https://arxiv.org/abs/2505.05295}.\par
NannyML documentation, \emph{Confusion matrix estimation} (v0.13.0), including confidence bands and alert thresholds.
\url{https://nannyml.readthedocs.io/en/v0.13.0/tutorials/performance_estimation/binary_performance_estimation/confusion_matrix_estimation.html}.\par
Białek, J., Kivimäki, J., Kuberski, W. and Perrakis, N. (2025; first posted 2024). Estimating model performance under covariate
shift without labels. arXiv:2401.08348v5.
\url{https://arxiv.org/abs/2401.08348v5}.\par
Angelopoulos, A.~N. and Bates, S. (2021; revised 2022). A gentle introduction to conformal prediction and distribution-free
uncertainty quantification. arXiv:2107.07511v6.
\url{https://arxiv.org/abs/2107.07511v6}.\par
NannyML documentation, \emph{Reverse concept drift},
\url{https://docs.nannyml.com/cloud/v0.22.0/model-monitoring/how-it-works/reverse-concept-drift-rcd}.\par

\srchead{Bias \& Fairness}
Angwin, J., Larson, J., Mattu, S. and Kirchner, L. (2016). Machine bias. \emph{ProPublica}.
\url{https://www.propublica.org/article/machine-bias-risk-assessments-in-criminal-sentencing}.\par
Barocas, S., Hardt, M. and Narayanan, A. (2023). \emph{Fairness and Machine Learning: Limitations and Opportunities}.
MIT Press. Chapter 3: Classification. \url{https://fairmlbook.org/classification.html}.\par
Chouldechova, A. (2017). Fair prediction with disparate impact: a study of bias in recidivism prediction instruments.
\emph{Big Data} 5(2). \url{https://arxiv.org/abs/1703.00056}.\par
Dieterich, W., Mendoza, C. and Brennan, T. (2016). \emph{COMPAS Risk Scales: Demonstrating Accuracy Equity and Predictive Parity}.
Northpointe Research Department.
\url{https://njoselson.github.io/pdfs/ProPublica_Commentary_Final_070616.pdf}.\par
\emph{Fairness and Algorithmic Decision Making}. Chapter 5: Parity measures.
\url{https://afraenkel.github.io/fairness-book/content/05-parity-measures.html}.\par
Google. Fairness: demographic parity. \emph{Machine Learning Crash Course}.
\url{https://developers.google.com/machine-learning/crash-course/fairness/demographic-parity}.\par
Hardt, M., Price, E. and Srebro, N. (2016). Equality of opportunity in supervised learning. \emph{NeurIPS}.
\url{https://arxiv.org/abs/1610.02413}.\par
Hu, X., Huang, Y., Li, B. and Lu, T. (2023). Inclusive FinTech lending via contrastive learning and domain adaptation.
arXiv preprint. \url{https://arxiv.org/abs/2305.05827}.\par
Kleinberg, J., Mullainathan, S. and Raghavan, M. (2017). Inherent trade-offs in the fair determination of risk scores.
\emph{ITCS}. \url{https://arxiv.org/abs/1609.05807}.\par
MLU-Explain. Equality of odds. \url{https://mlu-explain.github.io/equality-of-odds/}.\par
Verma, S. and Rubin, J. (2018). Fairness definitions explained. \emph{FairWare}.
\url{https://fairware.cs.umass.edu/papers/Verma.pdf}.\par

\srchead{Data Observability}
NIST/SEMATECH. What are control charts? \emph{e-Handbook of Statistical Methods}, Section 6.3.1.
\url{https://www.itl.nist.gov/div898/handbook/pmc/section3/pmc31.htm}.\par
NumPy. \texttt{numpy.quantile}. \emph{NumPy documentation}.
\url{https://numpy.org/doc/stable/reference/generated/numpy.quantile.html}.\par
PostgreSQL. Aggregate functions. \emph{PostgreSQL documentation}, Section 9.21.
\url{https://www.postgresql.org/docs/current/functions-aggregate.html}.\par
Soda. Metric monitoring dashboard. \emph{Soda documentation}.
\url{https://docs.soda.io/data-observability/metric-monitoring-dashboard}.\par
Soda. Unique count. \emph{Soda documentation}.
\url{https://docs.soda.io/data-observability/metric-monitoring-dashboard/column-monitors/all-data-types/unique-count}.\par
\endgroup
